\documentclass[11pt]{article}
\usepackage[margin=1.15in]{geometry}
\usepackage{amsmath,amssymb,amsthm}
\newtheorem{theorem}{Theorem}
\newtheorem{lemma}{Lemma}
\newtheorem{proposition}{Proposition}
\theoremstyle{remark}
\newtheorem{remark}{Remark}
\newcommand{\QQ}{\mathbb{Q}}
\DeclareMathOperator{\Lop}{L}

\title{Two closed forms for the Ap\'ery $\zeta(2)$ companion}
\author{}
\date{August 2026}
\begin{document}
\maketitle

\begin{abstract}
Let $S(n,k)=\binom nk^2\binom{n+k}k$ and let $(B_n)$ be the second solution of
Ap\'ery's $\zeta(2)$ recurrence. We prove
$5B_n=\sum_k S(n,k)\bigl(2\sum_{m\le n}\frac{(-1)^{m-1}}{m^2}
+\sum_{m\le k}\frac{(-1)^{n+m-1}}{m^2\binom nm\binom{n+m}m}\bigr)$
by an explicit certificate chain whose last step degenerates, remarkably, to the
identity $W(n,k)=3-\frac{3(n+1)}{n+1-k}$ together with the binomial theorem.
We also prove the depth-one harmonic (``minimal'') form by lifting its harmonic
letters to a two-parameter proper hypergeometric sum.  An order-three creative
telescope and an exact Ore-operator elimination reduce its mixed derivative to
Ap\'ery's order-two recurrence.  Every cell, boundary term, factorization, and
right remainder is checked in exact rational arithmetic.
\end{abstract}

\section{Statement}

Let
\[
S(n,k)=\binom nk^{2}\binom{n+k}k,\qquad a_n=\sum_{k=0}^{n}S(n,k),
\]
\[
(\Lop u)_n=(n+1)^{2}u_{n+1}-(11n^{2}+11n+3)u_n-n^{2}u_{n-1},
\]
so $\Lop a=0$, and let $(B_n)$ be the second solution: $\Lop B=0$, $B_0=0$,
$B_1=1$; classically $B_n/a_n\to\zeta(2)/5$. Put
\[
g_m(n)=\frac1{m^{2}\binom nm\binom{n+m}m}=\frac{((m-1)!)^{2}(n-m)!}{(n+m)!},
\qquad
\tau(n,k)=\sum_{m=1}^{\min(k,n)}(-1)^{n+m-1}g_m(n),
\]
\[
e(n)=\sum_{m=1}^{n}\frac{(-1)^{m-1}}{m^{2}},\qquad
w(n,k)=2e(n)+\tau(n,k).
\]

\begin{theorem}\label{thm:cl}
For all $n\ge0$,\quad $\displaystyle\sum_{k=0}^{n}S(n,k)\,w(n,k)=5B_n$.
\end{theorem}

This is the $\zeta(2)$ analogue of the classical representation of the
$\zeta(3)$ companion (van der Poorten \S8); we have not located it in the
literature in this form. The alternating factor $(-1)^{n}$ inside the inner sum
is essential.

\section{The miracle, and the two shifts of the tail}

Write $\sigma=(-1)^{n}$ and
$\gamma(n,k)=(-1)^{k}(k!)^{2}(n-k)!/(n+k+1)!$.

\begin{lemma}\label{lem:miracle}
For $0\le k\le n$,
\[
\tau(n+1,k)-\tau(n,k)=\frac{2\sigma}{n+1}\Bigl(\gamma(n,k)-\frac1{n+1}\Bigr),
\]
and for $0\le k\le n-1$,
$\ \tau(n-1,k)-\tau(n,k)=\frac{2\sigma}{n}\bigl(\gamma(n-1,k)-\frac1n\bigr)$.
\end{lemma}

\begin{proof}
$(-1)^{m}\bigl[g_m(n+1)+g_m(n)\bigr]
=2(n+1)\,(-1)^{m}\frac{((m-1)!)^{2}(n-m)!}{(n+m+1)!}$, while
\[
\gamma(n,m)-\gamma(n,m-1)
=(-1)^{m}\frac{((m-1)!)^{2}(n-m)!}{(n+m+1)!}\,
\bigl[m^{2}+(n+1-m)(n+1+m)\bigr],
\]
and $m^{2}+(n+1-m)(n+1+m)=(n+1)^{2}$. Summing over $m\le k$ telescopes;
$\gamma(n,0)=1/(n+1)$. The second identity is the first at $n-1$.
\end{proof}

Since $e(n\pm1)-e(n)$ cancels the constant terms,
\begin{equation}\label{eq:shifts}
w(n+1,k)=w(n,k)+\tfrac{2\sigma}{n+1}\gamma(n,k),\qquad
w(n-1,k)=w(n,k)+\tfrac{2\sigma}{n}\gamma(n-1,k).
\end{equation}

\section{Proof of Theorem \ref{thm:cl}}

Let $F(n,k)=S(n,k)w(n,k)$ and $\rho(n,k)=k^{2}+k(1+6n)-4-15n-11n^{2}$, and set
\[
G(n,k)=\frac{k^{3}\rho(n,k)}{n(n+1)^{2}}\binom{n+1}k^{2}\binom{n+k-1}k,
\qquad
U(n,k)=(-1)^{n+k}\binom nk .
\]

\emph{Step 1 (Zeilberger certificate).} For all $k\ge0$, $n\ge1$:
$(\Lop S(\cdot,k))_n=G(n,k+1)-G(n,k)$. After the absorptions
$\binom nk(n+1)=\binom{n+1}k(n+1-k)$ etc., this is a single polynomial identity
in $\QQ[n,k]$, verified by expansion; the shell form of $G$ makes it valid at
every cell, including $k=n,n+1$.

\emph{Step 2 (split).} By \eqref{eq:shifts} and the conversions
$S(n+1,k)\gamma(n,k)=(-1)^{k}\binom{n+1}k/(n+1-k)$,
$S(n-1,k)\gamma(n-1,k)=(-1)^{k}\binom{n-1}k/(n+k)$: for $0\le k\le n-1$,
\[
(\Lop F(\cdot,k))_n=(\Lop S(\cdot,k))_n\,w(n,k)
+U(n,k)\Bigl[\frac{2(n+1)^{2}}{(n+1-k)^{2}}-\frac{2(n-k)}{n+k}\Bigr].
\]

\emph{Step 3 (Abel).} Summing $k=0,\dots,n+1$, using Step 1, $G(n,0)=0$,
$G(n,n+2)=0$, $\Delta_k\tau=(-1)^{n+k}g_{k+1}$ and
$S(n,k+1)g_{k+1}(n)=\binom n{k+1}/(k+1)^{2}$, the sum of Step 2 collapses to
\[
(\Lop v)_n=\sum_{k=0}^{n}(-1)^{k}\binom nk\,W(n,k)\cdot(-1)^{n}
-2(n+1)^{2}+B_1(n)+B_2(n),
\]
where $v_n=\sum_kF(n,k)$,
$W=\frac{2(n+1)^{2}}{(n+1-k)^{2}}-\frac{2(n-k)}{n+k}
+\frac{k\rho(n,k)}{(n+1-k)^{2}(n+k)}$, and $B_1,B_2$ are the two boundary
cells $k=n,n+1$. Explicitly $B_1=2(n+1)^{2}$, and
$B_2=-1+2(n+1)\binom{2n+2}{n+1}(n!)^{2}/(2n+1)!=-1+4=3$.

\emph{Step 4 (the collapse).} As rational functions,
\[
\boxed{\;W(n,k)=3-\frac{3(n+1)}{n+1-k}\;}
\]
— both the double pole at $k=n+1$ and the pole at $k=-n$ cancel identically,
the former by the boundary identity $\rho(n,n+1)=-2(n+1)(2n+1)$. Hence, by
$\sum_k(-1)^{k}\binom nk=0$ $(n\ge1)$ and
\[
\sum_{k=0}^{n}\frac{(-1)^{k}\binom nk}{n+1-k}
\;\overset{k\mapsto n-k}{=}\;
(-1)^{n}\sum_{k=0}^{n}\frac{(-1)^{k}\binom nk}{k+1}
=\frac{(-1)^{n}}{n+1},
\]
(the last step by $\binom nk/(k+1)=\binom{n+1}{k+1}/(n+1)$ and the binomial
theorem), we get $\sum_k(-1)^{k}\binom nkW=3(-1)^{n+1}$ and therefore
$(\Lop v)_n=-3-2(n+1)^{2}+2(n+1)^{2}+3=0$ for $n\ge1$.

With $v_0=0$, $v_1=5=5B_1$ and $(n+1)^{2}\neq0$, induction gives $v_n=5B_n$. \qed

\section{The minimal form}

Let $w_{\min}=\tfrac15\bigl[H^{(2)}_n+H_k(2H_k-H_{n-k}-H_n)\bigr]$.

\begin{proposition}
$\sum_kS(n,k)w_{\min}(n,k)=B_n$ for all $n$ \emph{if and only if}
\begin{equation}\label{eq:L2}
(\Lop A)_n=-a_{n+1}-a_{n-1},\qquad
A_n:=\sum_{k=0}^{n}S(n,k)\,H_k\bigl(2H_k-H_{n-k}-H_n\bigr).
\end{equation}
\end{proposition}

\begin{proof}
$\Lop(H^{(2)}_\cdot a)_n=a_{n+1}+a_{n-1}$ exactly (the shifts of $H^{(2)}$
produce $1/(n+1)^{2}$, $-1/n^{2}$ against the leading coefficients), so
$\Lop(\sum_kSw_{\min})=\tfrac15[(a_{n+1}+a_{n-1})+(\Lop A)_n]$; initial values
$A_0=0$, $A_1=2$ match. \qed
\end{proof}

\begin{theorem}\label{thm:minimal}
For every $n\geq0$,
\[
 B_n=\frac15\sum_{k=0}^nS(n,k)
 \left[H_n^{(2)}+H_k(2H_k-H_{n-k}-H_n)\right].
\]
\end{theorem}

\begin{proof}
Put $X_m(t)=m!/(1-t)_m$ and
\[
 \Phi_n(u,v)=\sum_{k=0}^nS(n,k)X_k(u)
 \frac{X_k(v)^2}{X_{n-k}(v)X_n(v)}.
\]
At $(u,v)=(0,0)$ its coefficients $1,u,v,uv$ are respectively
$a_n,U_n,V_n,A_n$, where
\[
 U_n=\sum_kS(n,k)H_k,qquad
 V_n=\sum_kS(n,k)(2H_k-H_{n-k}-H_n).
\]
Let $E$ shift $n$ forward and set
\[
 \mathcal L=-(n+1)^2-(11n^2+33n+25)E+(n+2)^2E^2;
\]
thus $(\mathcal Lu)_n=(\Lop u)_{n+1}$.  Direct creative telescoping of the
proper hypergeometric summand produces an order-three operator
$\mathcal C(u,v)$ and a rational certificate.  Its denominator in $k$ is
$(n-k+1)^3(n-k+2)^3(n-k+3)^3$; the two endpoints and all ten top cells vanish
at the four Taylor coefficients above.

At the origin exact right factorizations give
\[
 \mathcal C(0,0)=Q_0\mathcal L,qquad
 \partial_v\mathcal C(0,0)=Q_v\mathcal L,
\]
where
\[
 Q_0=-2(2n+3)(29n+69)+(n+3)(29n+40)E
\]
and
\[
 Q_v=\frac{204n^3+1221n^2+2320n+1383}{(n+1)(n+3)}
 -\frac{138n^2+523n+468}{n+2}E.
\]
The $v$-derivative therefore gives $Q_0\mathcal LV=0$.  Since
$V_0=V_1=V_2=0$ and its leading coefficient never vanishes, $V=0$.

Write $C_u=\partial_u\mathcal C(0,0)$ and
$C_{uv}=\partial_u\partial_v\mathcal C(0,0)$.  Exact first-order operators
$X,Y$ satisfy $XQ_v=YQ_0$, and exact right division yields
\[
 YC_u-XC_{uv}-(XQ_0)J=H\mathcal L,qquad J=-(1+E^2),
\]
with zero remainder.  Differentiating the recurrence in $u$ and in $u,v$,
then eliminating $U$, gives
$XQ_0(\mathcal LA-Ja)=0$.  The leading coefficient of $XQ_0$ is
$(n+4)(29n+69)$; the first two values of $\mathcal LA-Ja$ vanish.  Hence
$\mathcal LA=Ja$, which is \eqref{eq:L2} after shifting.  The proposition now
proves the formula.
\end{proof}

\begin{remark}
The complete independently rerunnable audit is
\texttt{work/z2cf/parameter\_proof.mac}.  It prints the generic certificate
check, the four boundary cancellations, both factorizations, the common-left-
multiple identity, the zero Ore remainder, the nonzero leading coefficient,
and the two initial values.  No floating-point computation is used.
\end{remark}

\end{document}
